%% kapa: exam=1 round=29 year=2018 period=1 form=A subject=economics subject_ko=경제학원론 seq=11 number=51
%% source: Q-Net 감정평가사 기출문제 (cst00309) artlSeq=5202808 file=extracted/1008/5202808/2018년 제29회 감정평가사 1차시험 기출문제/2018년 제29회 감정평가사 1차시험 기출문제/1교시 A형.pdf
%% conversion: kordoc 4.24.51 -> tools/gen1.py (auto); review=proofread-sonnet; flags=none
\begin{question}{51}
복점(duopoly)시장에서 기업 A와 B는 각각 1, 2의 전략을 갖고 있다. 성과 보수 행렬(payoff matrix)이 다음과 같을 때, 내쉬균형의 보수쌍은? (단, 보수 행렬 내 괄호 안 왼쪽은 A, 오른쪽은 B의 보수이다.)
\par
{\small\centering
\begin{tabular}{|c|c|c|c|}\hline
\multicolumn{2}{|c|}{} & \multicolumn{2}{c|}{기업 B}\\\cline{3-4}
\multicolumn{2}{|c|}{} & 전략 1 & 전략 2\\\hline
\multirow{2}{*}{기업 A} & 전략 1 & (15,7) & (8,6)\\\cline{2-4}
 & 전략 2 & (3,11) & (10,7)\\\hline
\end{tabular}\par}
\begin{choices}
  \item (15,7)
  \item (8,6)
  \item (10,7)
  \item (3,11)과 (8,6)
  \item (15,7)과 (10,7)
\end{choices}
\end{question}
\answer{1}
